\scanpage{34}
\mspage{16}
Cela signifie qu'on se donne, pour $\alpha$ fixé,

1°) un graphe \inserted{connexe} $(S(\alpha),\widetilde A(\alpha),\deleted{\unsure{\widetilde A(\alpha)}}\sigma_{\widetilde A(\alpha)},\widetilde A(\alpha)\to S(\alpha))$

2°) on \deleted{\unsure{fix|se}} donne une application ``genre''
\[
S(\alpha)\xrightarrow{\gamma_\alpha}\mathbb N
\]

3°) on donne une application $I_\alpha\to S(\alpha)$, i.e. on ``répartit'' les \deleted{\unsure{divers}} ``pts marqués'' au dessus de $\alpha$ \inserted{(i.e. les éléments de $I_\alpha$)} entre les sommets $\beta\in S(\alpha)$ de $G(\alpha)$

4°) on donne une application $\widetilde A_\alpha\to S(\alpha)$, i.e. pour tt \deleted{flèche}\inserted{arc} du graphe $G$, d'origine $\alpha$, on \deleted{\unsure{se donne}} ``choisit une \inserted{nouvelle} origine'' parmi les sommets $\beta\in S(\alpha)$ du graphe $G(\alpha)$

5°) Les données 1°) : 4°) définissent une maquette, \inserted{0-connexe, $G(\alpha)$} et on prend soin que celle-ci soit un MD-graphe, \deleted{\unsure{ce}} ce qui signifie ici pour pour tt $\beta\in S(\alpha)$ ``de genre 0'' i.e. tel que $\gamma_\alpha(\beta)=0$, on ait $\operatorname{card}I_{\alpha,\beta}+\operatorname{card}\widetilde A_{\alpha,\beta}\geq 3$.

Ceci étant, on trouve un nouveau \inserted{MD-}graphe\inserted{$G'$}, \inserted{qui sera} spécialisé du \inserted{MD-}graphe $G$.

Je voudrais préciser de façon simplifiée la relation entre les MD-graphes $G$ et $G'$, et j'introduis à cet effet une digression sur
\mspage{33}
